\subsection{Приёмы отладки}
\subsubsection{Печать разреженной матрицы}
Для вывода разреженных матриц любого формата
(реализующего интерфейс \cvar{ISparseMatrix})
Можно воспользоваться функцией \cvar{dbg::print}
(её можно вызывать непосредственно в процессе отладки
не прерываясь на расстановку нужных вызовов в коде и перекомпиляцию).
\begin{cppcode}
#include "cfd/debug/printer.hpp"
...

LodMatrix M(3);
M.add_value(0, 0, 1.0); 
M.add_value(1, 0, 0.0); 
M.add_value(1, 1, 4.0); 
M.add_value(2, 2, 3.0); 
dbg::print(M);
\end{cppcode}
В результате в терминал выведется матрица
\begin{shelloutput}
-- SIZE = 3x3
     1      *      *
     0      4      *
     *      *      3
\end{shelloutput}
Здесь \cvar{*} -- отмечены места, где отсутсвует значение в шаблоне матрицы (то есть ноль по построению).
Заданные в шаблоне значения тоже могут быть нулями (как в настоящем примере на второй строке первой колонки).

Если матрица имеет большую размерность, то печатать всю матрицу целиком неудобно.
Но можно воспользоваться печатью отдельной строки.
Так вызов
\begin{cppcode}
dbg::print(1, M);
\end{cppcode}
выведет только значения второй строки матрицы, находящиеся в шаблоне
\begin{shelloutput}
-- ROW = 1
   [   0] 0   // M[1, 0] = 0
   [   1] 4   // M[1, 1] = 4
\end{shelloutput}

\subsubsection{Сохранение отладочных полей решения}
В процессе отладки работы с сеточными данными
часто возникает необходимость быстро отобразить
их на сетке. Для этого в файле \ename{debug/saver.hpp}
имеется набор вспомогательных функций (каждая из них может быть запущена непоследственно из отладчика):

\begin{cppcode}
// скаляр в узлах сетки
void save_point_data(const IGrid& grid, const std::vector<double>& data);
// скаляр в ячейках сетки
void save_cell_data(const IGrid& grid, const std::vector<double>& data);
// вектор в ячейках сетки
void save_cell_vector(const IGrid& grid, const std::vector<Vector>& vec);
// скаляр в расширенных точках коллокации (ячейки + граничные грани)
void save_extended_colloc_data(const IGrid& grid, const std::vector<double>& data);
// скаляр в гранях сетки
void save_face_data(const IGrid& grid, const std::vector<double>& data);
\end{cppcode}

Все эти функции сохраняют данные в файл \ename{dbg.vtk} в текущую папку из которой запущена программа.

\subsubsection{Профилирование}
\label{sec:dbg-tictoc}
Для профилирования (замера время исполнения участков кода) можно пользоваться
набором глобальных функций, определённых в файле
\ename{debug/tictoc.hpp}.

\paragraph{Простой замер времени}
Измерение времени исполнения участка кода:
\begin{cppcode}
#include "cfd/debug/tictoc.hpp"
...

dbg::Tic("process1");   // начать замер (блок с именем process1)
// ... код, время исполнение которого требуется замерить
std::this_thread::sleep_for(std::chrono::milliseconds(40));
dbg::Toc("process1");   // окончить замер
\end{cppcode}

По окончанию работы программы в терминал будет распечатан
отчёт о времени исполнения каждого блока (в миллисекундах)
\begin{shelloutput}
------ Tictoc report (ms)
Process        Total time    Calls count  Time per call
-------------------------------------------------------
process1               40              1             40
\end{shelloutput}

Названия секций может быть произвольным и может повторяться. В последнем случае
программа будет выводить как общее время исполнение этого блока, так и среднее время по запускам:

\begin{cppcode}
dbg::Tic("process1");
std::this_thread::sleep_for(std::chrono::milliseconds(40));
dbg::Toc("process1");

dbg::Tic("process1");
std::this_thread::sleep_for(std::chrono::milliseconds(50));
dbg::Toc("process1");
\end{cppcode}

\begin{shelloutput}
------ Tictoc report (ms)
Process        Total time    Calls count  Time per call
-------------------------------------------------------
process1               90              2             45
\end{shelloutput}

\paragraph{Замер времени внутри области видимости}
Чтобы не писать \cvar{Toc} для остановки таймера
можно использовать замер внутри текущей области видимомости (внутри текущих фигурных скобок) используя
метод \cvar{ScopedTic}.
Например, чтобы замерить время работы метода можно написать
\begin{cppcode}
void heavy_func(){
    auto tic = dbg::ScopedTic("heavy_func");
    std::this_thread::sleep_for(std::chrono::milliseconds(50));

    // объект tic удалится при выходы из области видимости
    // и таймер для процесса heavy_func будет остановлен
}
\end{cppcode}

\paragraph{Индексированные процессы}
Иногда бывает полезно разбить один процесс
на несколько подпроцессов. Для этого можно воспользоваться индексированными
именами. В отчёт при этом
пойдет как время исполнения каждого шага процесса, так и общее время.
\begin{cppcode}
{
    auto tic = dbg::ScopedTic("process", 1);
    std::this_thread::sleep_for(std::chrono::milliseconds(50));
}
{
    auto tic = dbg::ScopedTic("process", 2);
    std::this_thread::sleep_for(std::chrono::milliseconds(20));
}
{
    auto tic = dbg::ScopedTic("process", 3);
    std::this_thread::sleep_for(std::chrono::milliseconds(10));
}
\end{cppcode}
\begin{shelloutput}
------ Tictoc report (ms)
Process        Total time    Calls count  Time per call
-------------------------------------------------------
process                80              3             26
process-1              50              1             50
process-2              20              1             20
process-3              10              1             10
\end{shelloutput}

Это механизм удобно использовать для замера
времени исполнения тела цикла (если кол-во итераций небольшое)
\begin{cppcode}
for (size_t i=0; i<3; ++i){
    dbg::Tic("loop", i);
    std::this_thread::sleep_for(std::chrono::milliseconds(50-i));
    dbg::Toc("loop", i);
}
\end{cppcode}

\paragraph{Печать отчётов}
По умолчанию программа печатает
финальный отчёт о выполнении по окончании
работы программы.
Если требуется распечатать отчёт раньше, можно
явно вызывать метод \cvar{dbg::Report} или
\cvar{dbg::FinReport}.
Последний обнуляет все таймеры и может быть удобен
при запуске нескольких расчётов с разными параметрами:
\begin{cppcode}
for (size_t n_cells: {10, 20, 50}){
    calculate(n_cells);
    FinReport();  // печать отчёта о профилировании только для текущего расчёта
}
\end{cppcode}
